\documentclass[12pt]{article}
\usepackage[utf8]{inputenc}
\usepackage[T1]{fontenc}
\usepackage{hyperref}
\title{Central Atlas Tamazight Language Resources for K-12 Curriculum Development}
\author{Compiled by Claude Sonnet 5 (claude-sonnet-5) for Andr\'as Kornai}
\date{}
\begin{document}
\maketitle
\section{Introduction}
Central Atlas Tamazight (tzm) is spoken in the Middle Atlas and parts of the High Atlas in Morocco. In school it is taught as part of the standardized Moroccan Amazigh, not as a separate variety. This report rests on a small number of search results, mostly background pages.
\section{Primary Educational Resources}
IRCAM (Royal Institute of the Amazigh Culture) was created in 2001. Amazigh teaching began in primary schools in 2003. By 2008, IRCAM and the Ministry of Education had produced textbooks and teacher materials for all primary levels. A news item reports the presentation of the first Amazigh school textbook: \url{https://www.bladi.net/presentation-a-rabat-du-premier-manuel-scolaire-en-langue-amazighe.html}. A research article on Amazigh textbooks is at \url{https://l1research.org/article/download/117/115}. The search did not return the textbooks themselves.
\section{Linguistic and Reference Materials}
Standard Moroccan Amazigh blends Tashelhit, Central Atlas Tamazight and Tarifit, and has been an official language since 2011: \url{https://en.wikipedia.org/wiki/Standard_Moroccan_Amazigh}. A Leiden dissertation discusses standardization against the actual speech of Central Morocco, Tarifiyt and Tashelhiyt: \url{https://scholarlypublications.universiteitleiden.nl/handle/1887/4285960}. A French thesis (HAL) is at \url{https://halshs.archives-ouvertes.fr/tel-00669404v1/file/ThA_se_Vol1.pdf}; I did not verify its subject.
\section{Religious and Community Resources}
None were found in the search results.
\section{Digital and Audio Resources}
None were found in the search results.
\section{Curriculum Development Context}
The standardized Neo-Tifinagh-IRCAM script has been official since 2003. Existing school materials target the standard language, so a curriculum in Central Atlas Tamazight would have to adapt them to local speech. The sources found here are limited, and further searching is needed.
\end{document}